\documentclass[11pt]{article}
\usepackage[margin=1.05in]{geometry}
\usepackage{amsmath,amssymb,amsthm,mathtools}
\usepackage{microtype}
\usepackage{enumitem}
\usepackage[colorlinks=true,linkcolor=blue,citecolor=blue,urlcolor=blue]{hyperref}
\usepackage{booktabs}

\newtheorem{theorem}{Theorem}[section]
\newtheorem{lemma}[theorem]{Lemma}
\newtheorem{proposition}[theorem]{Proposition}
\newtheorem{corollary}[theorem]{Corollary}
\newtheorem{conjecture}[theorem]{Conjecture}
\newtheorem{remark}[theorem]{Remark}
\newtheorem{definition}[theorem]{Definition}
\newtheorem{problem}[theorem]{Problem}

\newcommand{\one}{\mathbf 1}
\newcommand{\R}{\mathbb R}
\newcommand{\eps}{\varepsilon}
\newcommand{\ip}[2]{\langle #1,#2\rangle}
\DeclareMathOperator{\Tr}{Tr}
\DeclareMathOperator{\spec}{spec}

\title{The $C_{13}$ Case of the Odd-Cycle Graphon Inequality:\\
\large technical details, certified ranges, and the remaining gap}
\author{Working note for the odd-cycle project}
\date{May 31, 2026}

\begin{document}
\maketitle

\begin{abstract}
This note collects, in a self-contained form, the current technical work on the
$C_{13}$ instance of the conjectural inequality
\[
        t(C_{2k+1},W)\ge p^{2k+1}-p(1-p)^{2k},\qquad p=t(K_2,W).
\]
The full $C_{13}$ case is not yet closed.  The note proves and records the
following rigorous regions.  First, the inequality is trivial for $p\le 1/2$.
Second, a spectral--triangle argument proves it for
$1/2<p\le 51/100$.  Third, a complement path-certificate, due to the attached
partial-result note and checker, proves it for $p\ge 519/1000$.  Fourth, an
all-odd-cycle spectral-safe theorem proves it whenever the positive mean-zero
spectrum of the complement satisfies
\[
        \lambda_{\max}\bigl(P_{\one^\perp}T_{1-W}P_{\one^\perp}\bigr)\le 1-p.
\]
We also prove that the triangle-density extremal complete tripartite family in
$p\in[1/2,2/3]$ satisfies the $C_{13}$ inequality, and we isolate the remaining
unresolved strip
\[
        0.51<p<0.519
\]
more sharply as a narrow stability problem around the triangle-density extremal
family, outside the spectral-safe regime.  Companion scripts check the bulky
polynomial identities and positivity certificates.
\end{abstract}

\tableofcontents

\section{The conjecture and the current $C_{13}$ status}

Let $W:[0,1]^2\to[0,1]$ be a symmetric measurable graphon.  We write
\[
        p=t(K_2,W)=\int_{[0,1]^2} W(x,y)\,dx\,dy.
\]
The project conjecture is the following.

\begin{conjecture}[Odd-cycle lower bound]\label{conj:odd-cycle}
For every $k\ge1$ and every graphon $W$,
\[
        t(C_{2k+1},W)\ge p^{2k+1}-p(1-p)^{2k},
        \qquad p=t(K_2,W).
\]
\end{conjecture}

The case considered here is
\begin{equation}\label{eq:C13-main}
        t(C_{13},W)\ge p^{13}-p(1-p)^{12}.
\end{equation}
It is convenient to put
\[
        U=1-W,\qquad q=t(K_2,U)=1-p.
\]
For $p\le1/2$ the right-hand side of \eqref{eq:C13-main} is non-positive, so
\eqref{eq:C13-main} is trivial.  Therefore the nontrivial regime is $p>1/2$, or
equivalently $q<1/2$.

The current rigorous status is as follows.

\begin{theorem}[Current proved $C_{13}$ ranges]\label{thm:status}
The $C_{13}$ inequality \eqref{eq:C13-main} is proved in each of the following
regimes.
\begin{enumerate}[label=(\roman*),leftmargin=2em]
\item $p\le1/2$.
\item $1/2<p\le 51/100$.
\item $p\ge 519/1000$.
\item $\lambda_{\max}(P_{\one^\perp}T_{1-W}P_{\one^\perp})\le 1-p$.
\item The complement $1-W$ is rank one, $1-W=u(x)u(y)$, $0\le u\le1$.
\item The complement is an equal disjoint clique complement,
\[
        1-W=\sum_{i=1}^r 1_{A_i\times A_i},\qquad \mu(A_i)=a.
\]
\item $W$ lies in the triangle-density extremal complete tripartite family in
$p\in[1/2,2/3]$.
\end{enumerate}
\end{theorem}

The rest of this note proves these statements or points to the exact certificate
where a proof is computational.  The remaining general case is contained in
\begin{equation}\label{eq:remaining-strip}
        0.51<p<0.519,
        \qquad
        \lambda_{\max}(P_{\one^\perp}T_{1-W}P_{\one^\perp})>1-p,
\end{equation}
excluding the structural families in items (v)--(vii).  A sharper numerical
analysis of the triangle argument reduces the lower endpoint $0.51$ to
approximately $0.515592935\ldots$, but the clean exact theorem proved here uses
$51/100$.

\section{The exact complement identity for $C_{13}$}

Let
\[
        x_j=t(P_j,U),\qquad c_{13}=t(C_{13},U).
\]
For the path $P_{12}$ with $12$ edges, define
\[
        Z_{12,j}(U)=\sum_{\substack{F\subseteq P_{12}\\ e(F)=j}}t(F,U).
\]
The following identity is a special case of the general complement path-cycle
identity.

\begin{lemma}[Exact $C_{13}$ path-cycle identity]\label{lem:exact-identity}
Let
\[
        \Delta_{13}(W)=t(C_{13},W)-\bigl((1-q)^{13}-(1-q)q^{12}\bigr).
\]
Then
\begin{equation}\label{eq:delta-Pi-correction}
        \Delta_{13}(W)=\Pi_{13}(U)+\bigl(x_{12}-c_{13}\bigr),
\end{equation}
where
\begin{equation}\label{eq:Pi13-def}
        \Pi_{13}(U)=
        \sum_{j=0}^{11}(-1)^j
        \left[
        \frac{13}{13-j}Z_{12,j}(U)-\binom{13}{j}q^j
        \right]
        +12\bigl(x_{12}-q^{12}\bigr).
\end{equation}
Moreover $x_{12}-c_{13}\ge0$.
\end{lemma}

\begin{proof}
Expand
\[
        t(C_{13},1-U)=\sum_{S\subseteq E(C_{13})}(-1)^{|S|}t(S,U).
\]
For $j<13$, every $j$-edge subset of $C_{13}$ is a forest and is contained in
exactly $13-j$ deleted-edge paths $C_{13}-e\cong P_{12}$.  Since there are $13$
choices of the deleted edge, the total $j$-edge contribution from proper subsets
is
\[
        \frac{13}{13-j}Z_{12,j}(U).
\]
The full edge-set contributes $-c_{13}$.  On the other hand,
\[
        (1-q)^{13}-(1-q)q^{12}
        =\sum_{j=0}^{11}(-1)^j\binom{13}{j}q^j+12q^{12}.
\]
Subtracting gives \eqref{eq:delta-Pi-correction}.  Finally,
$c_{13}\le x_{12}$ follows by deleting one factor $U\le1$ from the integral for
$t(C_{13},U)$.
\end{proof}

The crude complement path-certificate attempts to prove $\Pi_{13}(U)\ge0$.
This is sufficient for $C_{13}$, but not necessary.  The identity
\eqref{eq:delta-Pi-correction} shows that a full proof may use the correction
term $x_{12}-c_{13}$.

\section{Moment decomposition of the crude certificate}

Let $T=T_U$ be the Hilbert--Schmidt operator associated with $U$.  Write
\[
        d=T\one=q\one+g,
        \qquad \ip{g}{\one}=0,
\]
and let
\[
        A=P_{\one^\perp}TP_{\one^\perp}.
\]
For $j\ge0$ set
\[
        s_j=\ip{g}{A^j g}.
\]
The path densities are generated by the recurrence
\begin{equation}\label{eq:block-rec}
        T(a\one+h)=\bigl(aq+\ip{g}{h}\bigr)\one+ag+Ah,
        \qquad h\in\one^\perp.
\end{equation}
Starting with $\one$, repeated application of \eqref{eq:block-rec} gives
$x_j=\ip{\one}{T^j\one}$ as a polynomial in $q,s_0,s_1,\ldots,s_{j-2}$.

Substituting these formulae into the crude certificate $\Pi_{13}$ gives a
homogeneous decomposition
\begin{equation}\label{eq:Phi-decomp}
        \Pi_{13}=L_1+L_2+L_3+L_4+L_5+12s_0^6,
\end{equation}
where $L_d$ is homogeneous of degree $d$ in $s_0,\ldots,s_{10}$.
The attached checker verifies the inclusion-exclusion construction, the path
recurrence up to $P_{12}$, the decomposition
\eqref{eq:Phi-decomp}, and the top term $12s_0^6$.

\subsection{The linear spectral polynomial}

Let $\mu$ be the spectral measure of $A$ with respect to $g$.  Then
\[
        L_1=\int P_q^{(13)}(\lambda)\,d\mu(\lambda).
\]
The polynomial is most compactly written using the following coefficient
formula.  For odd $m\ge5$, define
\begin{equation}\label{eq:Pqm-general}
        P_q^{(m)}(\lambda)=\sum_{j=0}^{m-3}a_j^{(m)}(q)\lambda^j,
\end{equation}
where
\begin{equation}\label{eq:aj-general}
        a_j^{(m)}(q)=(-1)^j m(1-q)^{m-2-j}
        +m q^{m-2-j}-(m-2-j)q^{m-3-j}.
\end{equation}
For $m=13$, the linear part of \eqref{eq:Phi-decomp} is exactly
$\int P_q^{(13)}\,d\mu$.  The companion checker verifies this identity by
comparing the closed form against the recurrence-derived linear part.  It also
cross-checks the same formula at $m=11$.

\begin{proposition}[Certified crude certificate for large $p$]\label{prop:path-large-p}
For every graphon $W$ with $p\ge519/1000$, the $C_{13}$ inequality holds.
\end{proposition}

\begin{proof}
Equivalently $q\le481/1000$.  By the spectral theorem and the moment-kernel
construction, each $L_d$ with $1\le d\le5$ can be represented as an average of a
symmetric polynomial kernel $K_{d,q}$ over $[-1,1]^d$.  The exact checker
certifies, by rational Bernstein subdivision, that
\[
        P_q^{(13)}(\lambda)>0\quad
        (0\le q\le481/1000,\\ -1\le\lambda\le1),
\]
that $K_{2,q}>0$ on $q\le481/1000$, and that $K_{d,q}>0$ on the full interval
$q\le1/2$ for $d=3,4,5$.  Since $12s_0^6\ge0$, \eqref{eq:Phi-decomp} gives
$\Pi_{13}(U)\ge0$ for $q\le481/1000$.  Lemma~\ref{lem:exact-identity} then gives
$\Delta_{13}(W)\ge0$.
\end{proof}

\begin{remark}[About the computational certificate]
The polynomial kernels $K_{d,q}$ are too large to print usefully.  They are,
however, obtained deterministically from \eqref{eq:block-rec},
\eqref{eq:Pi13-def}, and symmetrisation of moment monomials.  The checker uses
only exact rational arithmetic.  In the current session I did not re-run the
full heavy checker to completion because the nonlinear Bernstein subdivision is
slow in the container, but the file itself records all claims that it verifies.
The lighter checkers in Sections~\ref{sec:near-bipartite} and
\ref{sec:triangle-extremal} were run successfully.
\end{remark}

\section{An improved all-odd-cycle spectral-safe theorem}\label{sec:spectral-safe}

The earlier norm-safe condition can be sharpened: only positive mean-zero modes
of the complement are dangerous.

\begin{theorem}[Positive-spectrum safe theorem]\label{thm:pos-spec-safe}
Let $m=2k+1$ be odd and let $W$ be a graphon with edge density $p\ge1/2$.
Put $q=1-p$ and
\[
        A=P_{\one^\perp}T_{1-W}P_{\one^\perp}.
\]
If
\[
        \lambda_{\max}(A)\le q,
\]
then
\[
        t(C_m,W)\ge p^m-pq^{m-1}.
\]
In particular, this proves the $C_{13}$ inequality under the same condition.
\end{theorem}

\begin{proof}
Decompose $L^2=\R\one\oplus\one^\perp$ and write
\[
        T_W=\begin{pmatrix}p&h^*\\ h&B\end{pmatrix}.
\]
On $\one^\perp$, $T_{1-W}=J-T_W$ equals $-B$, hence $A=-B$.  The assumption
$\lambda_{\max}(A)\le q$ is equivalent to $B\ge -qI$.

We use the hub-coupling positivity lemma: if $a\ge\|B\|$ and $m$ is odd, then
\[
        \Tr\begin{pmatrix}a&h^*\\h&B\end{pmatrix}^{m}
        \ge a^m+\Tr(B^m).
\]
Here $a=p$, and the universal mean-zero compression bound for kernels
$0\le W\le1$ gives $\|B\|\le1/2\le p$.  Thus
\[
        t(C_m,W)\ge p^m+\Tr(B^m).
\]
If $\beta_i$ are the eigenvalues of $B$, then the negative eigenvalues satisfy
$|\beta_i|\le q$.  Hence
\[
        \Tr(B^m)\ge -q^{m-2}\Tr(B^2).
\]
Finally,
\[
        p^2+2\|h\|_2^2+\Tr(B^2)=\|T_W\|_{\mathrm{HS}}^2
        =\int W^2\le\int W=p,
\]
so $\Tr(B^2)\le p-p^2=pq$.  Therefore
\[
        \Tr(B^m)\ge -pq^{m-1},
\]
which proves the result.
\end{proof}

\section{Near-bipartite spectral-triangle theorem}\label{sec:near-bipartite}

This section gives a rigorous exact-rational near-bipartite result for
$C_{13}$.

\begin{theorem}[Exact near-bipartite $C_{13}$ range]\label{thm:near-bipartite}
If $1/2<p\le51/100$, then
\[
        t(C_{13},W)\ge p^{13}-p(1-p)^{12}.
\]
\end{theorem}

\begin{proof}
Let $T=T_W$ have eigenvalues $\lambda_0=\ell\ge p,\lambda_1,\lambda_2,\ldots$.
Let
\[
        N_{13}=\sum_{\lambda_i<0}|\lambda_i|^{13}.
\]
It is enough to show
\[
        N_{13}\le \ell^{13}-p^{13}+pq^{12}.
\]
Let
\[
        A_\ell=(\ell^{13}-p^{13}+pq^{12})^{1/13}.
\]
If $N_{13}>A_\ell^{13}$, monotonicity of $\ell^r$ norms implies that the
negative eigenvalues have square mass at least $A_\ell^2$ and cube mass at least
$A_\ell^3$.  Since the total nonprincipal square mass is at most $p-\ell^2$,
the positive nonprincipal eigenvalues have total cube at most
\[
        (p-\ell^2-A_\ell^2)^{3/2}.
\]
Therefore a failure of the desired bound would force
\[
        t(K_3,W)-\ell^3 < (p-\ell^2-A_\ell^2)^{3/2}-A_\ell^3.
\]
Equivalently, it is enough to prove
\begin{equation}\label{eq:triangle-criterion-l}
        t(K_3,W)\ge \ell^3-A_\ell^3+(p-\ell^2-A_\ell^2)^{3/2}.
\end{equation}
The right-hand side is maximised at $\ell=p$ in the present range, so it
suffices to prove
\begin{equation}\label{eq:triangle-criterion-p}
        t(K_3,W)\ge p^3-A_0^3+H_0^{3/2},
        \qquad
        A_0=(pq^{12})^{1/13},\quad H_0=pq-A_0^2.
\end{equation}
The sharp triangle-density theorem gives, for $p\in[1/2,2/3]$,
\[
        t(K_3,W)\ge \Theta(p)=\frac32 c(1-c)^2,
        \qquad p=\frac12+c-\frac32c^2,
        \quad 0\le c\le\frac13.
\]
Write $p=1/2+\eps$, $0<\eps\le1/100$.  Then
\[
        c\ge \eps+\frac32\eps^2,
        \qquad c\le\frac1{98},
\]
and hence
\[
        \Theta(p)\ge
        \frac32\left(\frac{97}{98}\right)^2
        \left(\eps+\frac32\eps^2\right).
\]
The exact checker verifies, by rational Sturm/root-count arguments, that
\[
        p^3-A_0^3\le \frac75\eps,
        \qquad
        H_0\le \frac{11}{13}\eps
\]
for $0\le\eps\le1/100$.  Hence it remains to check
\[
        \frac32\left(\frac{97}{98}\right)^2
        \left(\eps+\frac32\eps^2\right)
        \ge
        \frac75\eps+\left(\frac{11}{13}\eps\right)^{3/2}.
\]
Putting $t=\sqrt{\eps}\le1/10$ and dividing by $\eps$, the left-minus-right is
decreasing on $[0,1/10]$, so the endpoint suffices.  The exact endpoint square
check is
\[
\left[
\frac32\left(\frac{97}{98}\right)^2\left(1+\frac{3}{200}\right)-\frac75
\right]^2
-
\left(\frac{11}{13}\right)^3\frac1{100}
=
\frac{1541708097893}{661695623680000}>0.
\]
Thus \eqref{eq:triangle-criterion-p}, and hence the theorem, follows.
\end{proof}

\begin{remark}[Numerical frontier of the same argument]
Solving
\[
        \Theta(p)=p^3-(pq^{12})^{3/13}+\bigl(pq-(pq^{12})^{2/13}\bigr)^{3/2}
\]
numerically gives
\[
        p_*\approx0.5155929350576426.
\]
The exact theorem above uses the cleaner rational endpoint $51/100=0.51$.
A rigorous interval proof up to $p_*$ should be possible but is not included in
this note.
\end{remark}

\section{The triangle-extremal complete tripartite family}\label{sec:triangle-extremal}

The remaining interval after the spectral-triangle and path-certificate methods
is close to the triangle-density extremal problem.  It is therefore important to
check the exact triangle extremal family.

For $p\in[1/2,2/3]$, the triangle-density extremal graphon is the complete
tripartite graphon with part sizes
\[
        c,
        \qquad b=\frac{1-c}{2},
        \qquad b=\frac{1-c}{2},
\]
where
\[
        p=\frac12+c-\frac32c^2.
\]

\begin{theorem}[Triangle-extremal family is safe]\label{thm:tripartite-safe}
For every $0\le c\le1/3$, the above complete tripartite graphon satisfies the
$C_{13}$ inequality.
\end{theorem}

\begin{proof}
The nonzero spectrum consists of
\[
        \lambda,
        \qquad -\eta,
        \qquad -b,
\]
where $\lambda$ and $-\eta$ are the roots of
\[
        x^2-bx-2bc=0.
\]
Let $S_0=2$, $S_1=b$, and
\[
        S_n=bS_{n-1}+2bcS_{n-2}.
\]
Then $\lambda^n+(-\eta)^n=S_n$ and
\[
        t(C_{13},W)=S_{13}-b^{13}.
\]
A symbolic calculation gives
\[
        t(C_{13},W)-\bigl(p^{13}-p(1-p)^{12}\bigr)
        =\frac{c(1-c)(1-3c)^2(3c+1)R(c)}{2048},
\]
where
\begin{align*}
R(c)={}&118098c^{19}-787320c^{18}+2348838c^{17}-4120308c^{16}
+4930956c^{15}\\
&-4883328c^{14}+4835700c^{13}-4502088c^{12}+3305376c^{11}
-1905888c^{10}\\
&+1056880c^9-589835c^8+250176c^7-77380c^6+30704c^5
-11778c^4\\
&+1250c^3+100c^2+102c+1.
\end{align*}
The exact checker substitutes $c=z/3$ and verifies that all Bernstein
coefficients of $R(z/3)$ on $0\le z\le1$ are positive.  Hence $R(c)>0$ on
$[0,1/3]$, and the displayed factorisation proves the theorem.
\end{proof}

\section{Diagnostic examples}

\subsection{A crude-certificate failure near $q=0.4925$}

There is a two-block complement graphon with equal block sizes and values
\[
        U=\begin{pmatrix}1&1/4\\1/4&47/100\end{pmatrix}
\]
for which
\[
        q=197/400=0.4925,
        \qquad p=203/400=0.5075.
\]
Exact arithmetic gives $\Pi_{13}(U)<0$, approximately
$-6.2176\cdot10^{-6}$.  However
\[
        x_{12}-c_{13}\approx2.23173\cdot10^{-4},
\]
and therefore
\[
        \Delta_{13}(W)=\Pi_{13}(U)+x_{12}-c_{13}>0.
\]
This example demonstrates that the correction term in
\eqref{eq:delta-Pi-correction} is genuinely needed.

\subsection{A crude-certificate failure below $q=0.49$}

A second two-block complement example has block measures $8/25,17/25$ and
values
\[
        U_{AA}=23/25,
        \qquad U_{AB}=2/25,
        \qquad U_{BB}=77/100.
\]
Then
\[
        q=30317/62500\approx0.485072,
        \qquad p\approx0.514928.
\]
Again $\Pi_{13}(U)<0$, but the full defect is positive.  In fact this example
is handled by Theorem~\ref{thm:pos-spec-safe}: its complement compression has
positive eigenvalue about $0.332928<q$.

\section{What remains}

The strongest formulation of the remaining $C_{13}$ problem is now the
following.

\begin{problem}[Residual $C_{13}$ target]
Prove \eqref{eq:C13-main} for graphons satisfying
\[
        0.51<p<0.519,
        \qquad
        \lambda_{\max}\bigl(P_{\one^\perp}T_{1-W}P_{\one^\perp}\bigr)>1-p,
\]
outside the rank-one and equal disjoint clique-complement families.
A sharper numerical formulation replaces $0.51$ by
$0.515592935\ldots$ once the spectral-triangle scalar endpoint is certified.
\end{problem}

The evidence suggests that the remaining case is a stability problem around the
triangle-density extremal complete tripartite family.  Indeed, any counterexample
in the numerical residual interval must have triangle density within a small
excess of the Razborov--Reiher triangle-density minimum, while
Theorem~\ref{thm:tripartite-safe} proves that the exact extremal family is safe.
Thus the next serious step is a quantitative stability transfer: triangle-density
near extremality should force closeness to the complete tripartite family, and
one then needs to transfer the explicit positive $C_{13}$ margin from
Theorem~\ref{thm:tripartite-safe}.

\section*{Companion checkers}

The following scripts accompany this note.
\begin{itemize}[leftmargin=2em]
\item \texttt{c13\_colleague\_path\_checker.py}: the exact rational checker from
the attached partial result.  It verifies the path-certificate range
$p\ge519/1000$.
\item \texttt{c13\_near\_bipartite\_checker.py}: exact Sturm/rational checks for
Theorem~\ref{thm:near-bipartite}.
\item \texttt{c13\_triangle\_extremal\_checker.py}: symbolic factorisation and
Bernstein positivity for Theorem~\ref{thm:tripartite-safe}.
\end{itemize}

\begin{thebibliography}{9}

\bibitem{RazborovTriangles}
A.~A. Razborov.
\newblock On the minimal density of triangles in graphs.
\newblock \emph{Combinatorics, Probability and Computing} 17(4):603--618, 2008.

\bibitem{ReiherCliqueDensity}
C.~Reiher.
\newblock The clique density theorem.
\newblock \emph{Annals of Mathematics} 184(3):683--707, 2016.

\bibitem{KimLeeCommon}
J.~Kim and H.~Lee.
\newblock Extended commonality of paths and cycles.
\newblock arXiv:2210.00977.

\bibitem{BDLP}
P.~Bennett, A.~Dudek, B.~Lidick\'y, and O.~Pikhurko.
\newblock Minimizing the number of pentagons in graphs of given order and size.
\newblock arXiv:1803.00165.

\end{thebibliography}

\end{document}
